% ============================================================================
% CHAPTER 3 - METHODS / PARAMETERS
% File: docs/thesis/methods/group3_connection_degradation.tex
% Group 3: grid connection, access road, transmission losses, degradation (Stage F6).
% Theory: docs/thesis/theory/ch2_grid_access_degradation.tex
% Values must match config/parameters.json. Lines marked % verify or % pending are open.
% ============================================================================

\section{Grid Connection and Degradation (Stage F6)}
\label{sec:params_group3}

This section defines how the cost of reaching the transmission network
and the decline of output over the plant life enter the LCOE of Stage F6.
It completes the cost structure of Section~\ref{sec:params_group2}.

\subsection{Point of Interconnection and Connection Length}
\label{subsec:g3_poi}

The point of interconnection of a cell is the nearest OpenStreetMap
substation whose nominal voltage is at least a country threshold
$V_{\mathrm{th},c}$, the lowest voltage at which a plant of about 100~MW
connects to a network able to absorb its output. In Brazil the threshold
is 138~kV, the subtransmission level at which plants of this size
commonly connect, below the 230~kV Basic Network. In India it is 220~kV,
the lowest voltage of the interstate transmission system to which
renewable plants of 50~MW or more may connect under the General Network
Access regulations \cite{cerc2022gna}. In Portugal it is 150~kV, the
lowest voltage of the transmission network operated by REN; the 60~kV
level belongs to the distribution network and has limited capacity for
plants of this size.
% verify: BRA 138 kV against ONS/ANEEL access rules (Procedimentos de Rede, PRODIST Modulo 3); IND 220 kV and the 50 MW limit in CERC GNA Regulations 2022
Substations without a voltage tag take the highest voltage of the OSM
lines that reach them; substations with no inferable voltage are
excluded, and their number is reported per country.

Distances are computed with the same geodesic distance transform used for
the other distance layers (Stage F2a), on the 0.01$^\circ$ grid, and
aggregated to each 0.05$^\circ$ cell as the mean over its eligible area.
The connection length is the distance corrected by a tortuosity factor:
\begin{equation}
    L_{\mathrm{grid},i} = \tau_{\mathrm{grid}}\, D_{\mathrm{sub},i},
    \qquad
    L_{\mathrm{road},i} = \tau_{\mathrm{road}}\, D_{\mathrm{road},i},
    \label{eq:g3_lengths}
\end{equation}
where $D_{\mathrm{sub},i}$ and $D_{\mathrm{road},i}$ are the distances to
the nearest eligible substation and to the nearest primary or secondary
road. Cells whose substation distance exceeds 100~km are flagged as
requiring transmission development rather than a radial spur; they remain
in the analysis with the same cost formula, and the share of potential
they hold is reported.

As a sensitivity, the substation distance is replaced by the distance to
the nearest line of the predictive network layer (gridfinder combined
with OSM \cite{arderne2020}), which has no voltage attribute and
therefore gives an optimistic, lower-bound connection length.

\subsection{Reference Project and Cost Boundary}
\label{subsec:g3_boundary}

Line, bay and road costs are incurred per project, not per megawatt.
They are converted to cost per unit capacity with a reference project of
$P_{\mathrm{ref}}=100$~MW for both technologies. This is a modelling
convention: it does not assume that every cell hosts a 100~MW plant, and
it is applied identically in all cells so that sites are compared on the
same basis.

The capital cost of Section~\ref{subsec:g2_capex} covers everything
inside the plant boundary: collection networks, collector substations,
step-up transformers and the switchyard. For PV this matches the scope
stated by IRENA \cite{irena2025costs2024}. For onshore wind the report
does not state its scope, and the same boundary is assumed; this is
recorded as a limitation. Outside the boundary GeoFREA adds the external
interconnection cost per unit capacity
\begin{equation}
    C_{\mathrm{inter},i} =
    \frac{c_{\mathrm{line}}\,L_{\mathrm{grid},i}
          + C_{\mathrm{bay}}
          + c_{\mathrm{road}}\,L_{\mathrm{road},i}}{P_{\mathrm{ref}}}
    \quad [\mathrm{USD/MW}],
    \label{eq:g3_inter}
\end{equation}
with $c_{\mathrm{line}}$ the cost of a single-circuit overhead line at
$V_{\mathrm{th},c}$ (USD/km), $C_{\mathrm{bay}}$ the cost of the
connection works at the receiving substation (USD) and
$c_{\mathrm{road}}$ the cost of a new access road (USD/km). Upstream
network reinforcement is not included, so the result is the cost of the
radial connection only, not a full interconnection-study estimate
\cite{seel2023}.

Unit costs come from regulatory or regulator-compiled benchmarks of each
country (Table~\ref{tab:g3_costs}). For India, the line cost is the
indicative capital cost of a 220~kV single-circuit line published by the
Central Transmission Utility, which excludes right of way
\cite{ctuilIndicative2025}. For Portugal, no national per-kilometre
benchmark is public; the line cost is the ACER indicator for 220~kV
single-circuit overhead lines, which is conservative for 150~kV, and the
connection works are represented by the ACER indicator for the upgrade of
an existing air-insulated substation with up to five bays
\cite{acer2023uic}. For Brazil, line and bay costs come from the ANEEL
reference price bank \cite{aneelBPR}. All values are converted to 2024
USD with the national GDP deflator and the 2024 average exchange rate.
Each cost is sampled uniformly between the bounds in the table.

\begin{table}[htbp]
    \centering
    \caption{Unit costs of the external connection, in source currency and price year.}
    \label{tab:g3_costs}
    \small
    \begin{tabularx}{\textwidth}{@{}l
        >{\raggedright\arraybackslash}X
        >{\raggedright\arraybackslash}p{3.3cm}
        >{\raggedright\arraybackslash}p{2.3cm}@{}}
        \toprule
        Country & Item & Central [range] & Source \\
        \midrule
        IND & $c_{\mathrm{line}}$, 220~kV S/C, single Moose / AL59 Zebra
            & 1.10 [1.10, 1.14] INR crore/km (2025)
            & \cite{ctuilIndicative2025} \\
        IND & $C_{\mathrm{bay}}$, 220~kV line bay
            & to be sourced & -- \\
        PRT & $c_{\mathrm{line}}$, 220~kV S/C (used at 150~kV)
            & 0.362 [0.303, 0.543] MEUR/km
            & \cite{acer2023uic} \\
        PRT & $C_{\mathrm{bay}}$, AIS substation upgrade, 0--5 bays
            & 2.789 [2.116, 4.300] MEUR
            & \cite{acer2023uic} \\
        BRA & $c_{\mathrm{line}}$, 138~kV S/C
            & to be sourced & \cite{aneelBPR} \\
        BRA & $C_{\mathrm{bay}}$, 138~kV line entry module
            & to be sourced & \cite{aneelBPR} \\
        All & $c_{\mathrm{road}}$, unpaved access road
            & to be sourced & -- \\
        \bottomrule
    \end{tabularx}
\end{table}
% verify: ACER (2023) price base year; ranges for PRT are the interquartile range (Q1-Q3), central is the median
% pending: IND bay, BRA line and bay, road costs

\subsection{Transmission Losses}
\label{subsec:g3_losses}

The energy of Stage F5 is measured at the plant terminals. Resistive
losses on the spur line are not modelled in the central case, because no
country-specific line parameters were found; the energy is used as
delivered at the POI ($\eta_{\mathrm{trans}}=1$). As a sensitivity, a
linear loss of $\alpha_{\mathrm{loss}}$ per kilometre is applied,
\begin{equation}
    \eta_{\mathrm{trans},i} = 1 - \alpha_{\mathrm{loss}}\,L_{\mathrm{grid},i},
    \qquad \alpha_{\mathrm{loss}} \in [0.03,\,0.05]~\%/\mathrm{km},
    \label{eq:g3_losses}
\end{equation}
which is author judgment (Tier 3). For a 100~MW plant connected at
138~kV, the resistive loss at rated flow is of the order of
$P\,r'\,L/V^2$, a few percent at 50~km for typical conductor resistance,
which is consistent with that range.

\subsection{Degradation}
\label{subsec:g3_degradation}

Annual energy declines geometrically,
\begin{equation}
    E_{i,t} = E_{i,1}\,(1-d)^{t-1}, \qquad t = 1,\dots,n,
    \label{eq:g3_degradation}
\end{equation}
which is the form already used in the Stage F6 kernel.

For PV, $d$ follows the temperature zone of the cell. Zones are taken
from the global PV climate-zone dataset of Karin et al.\ \cite{karin2019,pvcz},
the same scheme used to classify the fleet of Jordan et al.\
\cite{jordan2022}: the zone of a cell is that of the nearest 0.25$^\circ$
point of the dataset, using the equivalent temperature of an open-rack
module. The zone is a property of the site's historical climate and does
not change across climate members. The rate of each zone is computed
from the public fleet data of Jordan et al.\ \cite{jordan2022}: the
central value is the zone median and the sampling range runs from the
median to the mean, so the uncertainty is skewed toward faster
degradation, as in the fleet data. The fleet covers zones T1 to T7;
cells in hotter zones take the T7 value, and zones with too few systems
are merged with the adjacent zone. Most of the territory of Brazil and
India lies in zones T6 and T7 and most of Portugal in T5, so the hot end
of the fleet sample determines the result.
% pending: zone rates computed from Fleet_results_public_2022_02_02.csv (DuraMAT); minimum systems per zone set in config

For onshore wind, $d$ follows the vintage of the reference turbine. The
central value is the rate of newer vintages and the sampling range
extends to that of older vintages \cite{hamilton2020}. The policy-related
drop observed after ten years in the United States is not included. The
complex-terrain rate of Kim and Kim \cite{kim2025} is kept as a
deterministic sensitivity, not sampled.

\begin{table}[htbp]
    \centering
    \caption{Annual degradation rate (uniform sampling).}
    \label{tab:g3_degradation}
    \small
    \begin{tabular}{llccc}
        \toprule
        Technology & Class & Central & Range & Source \\
        \midrule
        PV   & zones T1--T7 & zone median & median--mean & \cite{jordan2022} \\
        Wind & reference turbine & 0.17\,\%/yr & 0.17--0.53\,\%/yr & \cite{hamilton2020} \\
        Wind & complex terrain (sensitivity) & 0.72\,\%/yr & -- & \cite{kim2025} \\
        \bottomrule
    \end{tabular}
\end{table}

All values are transferred from United States fleets (Tier 2). The fleet
medians of about 0.48\,\%/yr in cooler and 0.88\,\%/yr in hotter climates
\cite{jordan2022} serve as a check on the zone values, and the global
meta-analysis median of 0.94\,\%/yr \cite{straubmuck2025} is reported for
comparison only.

\subsection{Levelized Cost with Connection and Degradation}
\label{subsec:g3_lcoe}

With these terms, the LCOE of cell $i$ in climate member $m$ and
parameter draw $s$ is
\begin{equation}
    \mathrm{LCOE}_{i,m,s} =
    \frac{I_s + C_{\mathrm{inter},i,s}
          + \sum_{t=1}^{n_s} F_s\,(1+r_s)^{-t}}
         {\sum_{t=1}^{n_s} E_{i,m,1}\,(1-d_{i,s})^{t-1}\,
          \eta_{\mathrm{trans},i}\,(1+r_s)^{-t}},
    \label{eq:g3_lcoe}
\end{equation}
where $I$ is the plant capital cost, $F=f\,I$ the fixed operating cost
of Section~\ref{subsec:g2_opex}, $r$ the real discount rate and $n$ the
lifetime. The connection cost enters as capital at commissioning. A cell
with a missing distance is not given a zero connection cost: it is
flagged and excluded from the LCOE results, and the count is reported.

\subsection{Parameters}
\label{subsec:g3_parameters}

\begin{table}[htbp]
    \centering
    \caption{Group 3 parameters.}
    \label{tab:g3_summary}
    \small
    \begin{tabularx}{\textwidth}{@{}>{\raggedright\arraybackslash}p{3.6cm}
        >{\raggedright\arraybackslash}X
        >{\centering\arraybackslash}p{1.3cm}
        >{\centering\arraybackslash}p{0.9cm}@{}}
        \toprule
        Parameter & Value & Sampled & Tier \\
        \midrule
        $V_{\mathrm{th}}$ & BRA 138~kV; IND 220~kV; PRT 150~kV & no & 2 \\
        $P_{\mathrm{ref}}$ & 100~MW & no & 3 \\
        $\tau_{\mathrm{grid}}$ & 1.25, U[1.2, 1.3] & yes & 3 \\
        $\tau_{\mathrm{road}}$ & 1.25, U[1.2, 1.3] & yes & 3 \\
        $c_{\mathrm{line}}$, $C_{\mathrm{bay}}$, $c_{\mathrm{road}}$ & Table~\ref{tab:g3_costs} & yes & 2 \\
        $\alpha_{\mathrm{loss}}$ & 0 (central); 0.03--0.05\,\%/km (sensitivity) & no & 3 \\
        $d$ & Table~\ref{tab:g3_degradation} & yes & 2 \\
        Long-distance flag & 100~km & no & 3 \\
        \bottomrule
    \end{tabularx}
\end{table}

\subsection{Limitations}
\label{subsec:g3_limitations}

The connection cost covers the radial spur, the works at the receiving
substation and an access road, not network reinforcement, which can
dominate real interconnection costs. Substation voltage comes from OSM
tags, which are incomplete, most of all in India. Tortuosity is a uniform
factor rather than a routed path. Line costs exclude right of way in
India and come from a higher voltage class in Portugal. The capital-cost
boundary of IRENA is stated for PV and assumed for onshore wind.
Degradation rates come from United States fleets, whose hottest zone is
the most common one in Brazil and India; the PV zone is fixed at the
historical climate, and wind uses a single vintage. Transmission losses
are not modelled in the central case.
